%%%%%%%%%%%%%%%%%%%%%%%%%%%%%%%%%%%%%%%%%%%%%%%%%%%%%%%%%%%%
%  Local Income Taxes and Residential Sorting --- Pitch
%%%%%%%%%%%%%%%%%%%%%%%%%%%%%%%%%%%%%%%%%%%%%%%%%%%%%%%%%%%%

\documentclass[xcolor = x11names, compress, aspectratio = 169]{beamer}
\setbeamercovered{transparent}

\usepackage{graphicx}
\usepackage{booktabs}
\usepackage{hyperref}

\usefonttheme[onlymath]{serif}
\beamertemplatenavigationsymbolsempty

\xdefinecolor{GopherMaroon}{HTML}{7A0019}
\xdefinecolor{GopherGold}{HTML}{FFCC33}
\xdefinecolor{GopherDarkMaroon}{HTML}{5B0013}

\useoutertheme[subsection = false, shadow]{miniframes}
\useinnertheme{default}
\usefonttheme{structureitalicserif}
\usepackage{palatino}
\setbeamerfont{title like}{shape = \scshape}
\setbeamerfont{frametitle}{shape = \scshape}

\setbeamercolor*{lower separation line head}{bg = GopherMaroon}
\setbeamercolor*{normal text}{fg = black, bg = white}
\setbeamercolor*{alerted text}{fg = GopherDarkMaroon}
\setbeamercolor*{structure}{fg = GopherDarkMaroon}

\setbeamertemplate{footline}
{
    \leavevmode
    \hbox{%
        \begin{beamercolorbox}[wd = .333333\paperwidth, ht = 2.25ex, dp = 1ex, center]{}
            Hockert
        \end{beamercolorbox}
        \begin{beamercolorbox}[wd = .333333\paperwidth, ht = 2.25ex, dp = 1ex, center]{}
            Local Income Taxes and Residential Sorting
        \end{beamercolorbox}
        \begin{beamercolorbox}[wd = .333333\paperwidth, ht = 2.25ex, dp = 1ex, right]{}
            \today \hspace*{2em}
            \insertframenumber{} / \inserttotalframenumber\hspace*{2ex}
        \end{beamercolorbox}}
        \vskip0pt
    }

\hypersetup{colorlinks, linkcolor = black, urlcolor = GopherMaroon}

\newcommand{\src}[2]{\textsuperscript{\href{#1}{\footnotesize[#2]}}}

\begin{document}

\begin{frame}
    \fontfamily{cmr}
    \title{Local Income Taxes and Residential Sorting \\
    \large Pitch}
    \author[Hockert]{Matthew Hockert}
    \date{\small{\textit{\today}}}
    \institute{\textit{Applied Economics, University of Minnesota}}
    \titlepage
\end{frame}

%======================================================================
\begin{frame}
\fontfamily{cmr}
\frametitle{The question}
    \begin{itemize}\setlength{\itemsep}{0.3cm}
        \item 17 states let local governments tax earned
        income\src{https://taxfoundation.org/research/all/state/local-income-taxes-2019/}{TF}.
        \item Do people move to avoid income taxes or do they take on the cost?
        \item Can't find research specifically answering this question at the local level. 
        \item Specifically looking at sub-county where standard migration data (ACS, IRS SOI) is too coarse. (most of studied state level policies)
    \end{itemize}
\end{frame}

%======================================================================

\begin{frame}
\fontfamily{cmr}
\frametitle{Two margins --- both need more research}
    \begin{itemize}\setlength{\itemsep}{0.3cm}
        \item Rate changes (intensive): do higher rates push workers
        out? PA, MD
        \src{https://www.marylandtaxes.gov/individual/income/tax-info/tax-rates.php}{MD}, MI, and KY
        \src{https://www.klc.org/InfoCentral/Detail/31/occupational-license-tax}{KY}.
        \item Adoption (extensive): does turning a tax on cause
        sorting?:
        \begin{itemize}\setlength{\itemsep}{0.05cm}
            \item School Districts
            \item Cities
            \item Counties
        \end{itemize}
        \item Endogeneity of ``why are you raising (or lowering?) income taxes?''
    \end{itemize}
\end{frame}


\begin{frame}
\fontfamily{cmr}
\frametitle{Why I want to start with Pennsylvania (intensive margin)}
    \begin{itemize}\setlength{\itemsep}{0.3cm}
        \item Every city and school district levies an earned income
        tax which gives continuous variation on rates over geographies
        \src{https://dced.pa.gov/local-government/local-income-tax-information/psd-codes-and-eit-rates/}{rates}.
        \item Every city files an Annual Financial Report so we can
        see what the money funds and what other taxes were raised or lowered
        \src{https://apps.dced.pa.gov/munstats-public/ReportInformation2.aspx?report=mAfrForm}{AFR}. (can directly address endogeneity concerns). Can easily make a scraper for this site.
        \item Act 32 (2012) centralized collection and made employer
        withholding binding a natural experiment for whether
        enforcement matters
        \src{https://business.pa.gov/operate/hiring-workers/act-32-and-local-earned-income-tax/}{Act 32}.
    \end{itemize}
\end{frame}

%======================================================================
\begin{frame}
\fontfamily{cmr}
\frametitle{Why LODES (origin--destination data)}
    \begin{itemize}\setlength{\itemsep}{0.3cm}
        \item Links where you live to where you work, annually, at
        the Census block
        \src{https://lehd.ces.census.gov/data/lodes/LODES8/}{LODES}.
        \item I work with LODES data for my RA at Accessibility Observatory (familiarity)
        \item Use aggregate data at Census Tract to look at shifts within cities/school districts that the tract is in.
        \item When a city raises its rate, does the
        share of its workers who also live there fall?
    \end{itemize}
\end{frame}

%======================================================================
\begin{frame}
\fontfamily{cmr}
\frametitle{How PA tax works (post-Act 32, ex.\ Philadelphia)}
    \small
    The worker pays the higher of the home resident rate or the work
    nonresident rate\src{https://dced.pa.gov/local-government/local-income-tax-information/local-withholding-tax-faqs/}{DCED}, but how that bill is split between the cities is
    piecewise:
    \[
        \bigl(\tau^{\text{home}},\ \tau^{\text{work city}}\bigr) \;=\;
        \begin{cases}
            \bigl(\tau^{\text{res}} - \tau^{\text{work,nr}},\ \tau^{\text{work,nr}}\bigr)
                & \text{if } \tau^{\text{res}} > \tau^{\text{work,nr}} \\[2pt]
            \bigl(0,\ \tau^{\text{work,nr}}\bigr)
                & \text{if } \tau^{\text{res}} \leq \tau^{\text{work,nr}}
        \end{cases}
    \]
    \begin{itemize}\setlength{\itemsep}{0.18cm}
        \item Work city always collects its full nonresident rate.
        \item Home city only collects the leftover, and gets nothing
        if the work rate is higher.
        \item Which rate is higher also determines the active margin:
        home-binds $\rightarrow$ move home; work-binds $\rightarrow$
        change jobs (or absorb the tax).
        \item Philadelphia operates under the Sterling Act and is
        excluded from Act 32's credit/max
        system\src{https://dced.pa.gov/local-government/local-income-tax-information/act32-faq/}{Act 32 FAQ}.
        Withholding only became binding statewide on
        Jan.\ 1, 2012\src{https://dced.pa.gov/local-government/local-income-tax-information/act32-faq/}{DCED}.
    \end{itemize}
\end{frame}

%======================================================================
\begin{frame}
\fontfamily{cmr}
\frametitle{Example: a \$50{,}000 worker}
    \small
    \begin{itemize}\setlength{\itemsep}{0.18cm}
        \item Case 1: lives in $B$ (1.0\%), works in $A$ (0.5\%).
        Bill $=$ \$500. If $B$ raises to 2.0\%, bill $=$ \$1{,}000.
        Home rate binds, so she escapes only by moving her home.
        Expect $B$'s resident-worker count to fall.
        \item Case 2: lives in $A$ (0.5\%), works in $B$ (1.0\%).
        Bill $=$ \$500. If $B$ raises to 2.0\%, bill $=$ \$1{,}000.
        Work rate binds, so she escapes only by changing jobs.
        Expect $B$'s total-jobs count to fall.
        \item LODES records each worker as a home-tract to work-tract
        flow, so either kind of move shows up directly in the data.
    \end{itemize}
\end{frame}

%======================================================================
\begin{frame}
\fontfamily{cmr}
\frametitle{Methods (proposed)}
    \small
    Staggered event study on a muni $\times$ year panel:
    \[
        Y_{it} = \alpha_i + \gamma_t
                 + \sum_{\ell} \mu_\ell \cdot D_{it}^{\ell}
                 + \varepsilon_{it}
    \]
    \begin{itemize}\setlength{\itemsep}{0.2cm}
        \item Outcomes of interest:
        \begin{itemize}\setlength{\itemsep}{0.04cm}
            \item Share of a city's workers who also live there
            \item Resident workers, nonresident workers, total jobs
            \item Fiscal response (EIT revenue, spending by function)
            \item Muni rate vs.\ school district rate
        \end{itemize}
        \item Robustness: de Chaisemartin \& D'Haultfoeuille (2024)
        continuous staggered estimator, which handles rate size
        rather than just on/off
        \src{https://doi.org/10.3982/ECTA19039}{DCDH 2024}.
    \end{itemize}
\end{frame}

\end{document}
